\documentclass[11pt,a4paper]{article}
\usepackage[T1]{fontenc}
\usepackage[utf8]{inputenc}
\usepackage{lmodern}
\usepackage{amsmath,amssymb,amsthm,mathtools}
\usepackage{esint,graphicx}
\usepackage[a4paper,margin=22mm,headheight=14pt]{geometry}
\usepackage{microtype}
\usepackage{booktabs,longtable,array,enumitem}
\newcolumntype{L}[1]{>{\raggedright\arraybackslash}p{#1}}
\usepackage[dvipsnames]{xcolor}
\usepackage{fancyhdr}
\usepackage{hyperref}

\newtheorem{theorem}{Theorem}[section]
\newtheorem{lemma}[theorem]{Lemma}
\newtheorem{proposition}[theorem]{Proposition}
\newtheorem{corollary}[theorem]{Corollary}
\theoremstyle{remark}
\newtheorem{remark}[theorem]{Remark}
\newtheorem{definition}[theorem]{Definition}
\newtheorem{claim}[theorem]{Claim}
\setlist{nosep,leftmargin=1.5em}
\setlength{\emergencystretch}{2em}
\pagestyle{fancy}
\fancyhf{}
\fancyhead[L]{\small AIM 689}
\fancyhead[R]{\small Revised 7 November 2026}
\fancyfoot[C]{\thepage}
\renewcommand{\headrulewidth}{0.3pt}
\hypersetup{colorlinks=true,linkcolor=MidnightBlue,urlcolor=MidnightBlue,citecolor=MidnightBlue,pdftitle={AIM 689},pdfauthor={Matthew J. Colbrook; results found by AI}}
\begin{document}
\title{Failure of the sharp affine-modulus bound}
\author{Matthew J. Colbrook\thanks{The solutions and partial results in this document were found by AI. The supplied research and the separate repository review are AI work; human peer review and formal verification are not claimed.}\\[4pt]
\small Department of Applied Mathematics and Theoretical Physics\\
\small University of Cambridge, Cambridge CB3 0WA, United Kingdom\\
\small\href{mailto:m.colbrook@damtp.cam.ac.uk}{m.colbrook@damtp.cam.ac.uk}}
\date{7 November 2026}
\maketitle
\begin{abstract}
A small explicit harmonic perturbation of the critical circular-annulus map violates the proposed bound even after optimizing over all real affine transformations. The conclusion concerns exactly the catalogue target; no historical priority claim is made.
\end{abstract}
\noindent\textbf{Catalogue scope.} AIM 689; complete resolution. \href{https://github.com/MColbrook/AIM/blob/1e564633494c109edfd3ec46da1366adcff925a5/problems/689-a-sharp-affine-modulus-bound-for-harmonic-homeomorphisms.md}{Canonical target at revision \texttt{1e56463}}.


\section{Statement and argument}
\label{sec:689}

The sharp necessary lower bound proposed in \cite[Conjecture 1.3]{IwaniecKovalevOnninen} fails under a small explicit harmonic perturbation of the critical circular-annulus map. The affine optimization is included below; decreasing ordinary modulus alone would not settle the question.

\begin{theorem}
Fix $\tau>0$ and put $R=\cosh\tau$. For every sufficiently small nonzero real $\varepsilon$, the harmonic map
\begin{equation}
 h_\varepsilon(z)=\frac12\left(z+\frac1{\bar z}\right)
       +\frac\varepsilon2\left(z^2+\frac1{\bar z^2}\right),
 \qquad 1<|z|<e^\tau,
 \label{eq:689-map}
\end{equation}
is a homeomorphism onto a bounded ring $\Omega_\varepsilon$ with analytic boundary. Its target satisfies
\[
 \operatorname{Mod}_{\mathrm{aff}}(\Omega_\varepsilon)
 <\log R=\log\cosh\!\left(\operatorname{Mod}\{1<|z|<e^\tau\}\right).
\]
Thus the inequality in Problem 689 is false.
\end{theorem}

\begin{proof}
Write
\[
 S=\cosh(2\tau)=2R^2-1,\qquad m=\log R>0.
\]
In logarithmic polar coordinates,
\begin{equation}
 h_\varepsilon(e^{t+i\theta})
 =\cosh t\,e^{i\theta}+\varepsilon\cosh(2t)\,e^{2i\theta}.
 \label{eq:689-polar-map}
\end{equation}
Its two boundary curves are
\begin{equation}
 \Gamma_{0,\varepsilon}(\theta)=e^{i\theta}+\varepsilon e^{2i\theta},
 \qquad
 \Gamma_{1,\varepsilon}(\theta)=R e^{i\theta}+\varepsilon S e^{2i\theta}.
 \label{eq:689-boundaries}
\end{equation}
For sufficiently small $|\varepsilon|$, these are analytic Jordan curves. They are nested, for example, whenever $|\varepsilon|<(R-1)/(S+1)$, since the inner curve lies in the disk of radius $1+|\varepsilon|$ while the outer curve encloses the disk of radius $R-|\varepsilon|S$. Let $\Omega_\varepsilon$ be the ring between them.

\paragraph{Homeomorphism.}
The Jacobian of \eqref{eq:689-polar-map} with respect to $(t,\theta)$ is
\begin{equation}
 J_\varepsilon(t,\theta)
 =\frac12\sinh(2t)+2\varepsilon\sinh(3t)\cos\theta
       +2\varepsilon^2\sinh(4t).
 \label{eq:689-jacobian}
\end{equation}
Because $\sinh(3t)/\sinh(2t)$ extends continuously to $t=0$ and is bounded on $[0,\tau]$, this is positive for $0<t\le\tau$ when $|\varepsilon|$ is small enough. Each intermediate curve is a Jordan curve: the polynomial
\[
 w\longmapsto\cosh t\,w+\varepsilon\cosh(2t)w^2
\]
is injective on the closed unit disk if $2|\varepsilon|\cosh(2t)<\cosh t$, which holds uniformly for $0\le t\le\tau$ after reducing $|\varepsilon|$. The boundary winding degree is one between \eqref{eq:689-boundaries}, and zero outside. Every interior preimage has positive Jacobian by \eqref{eq:689-jacobian}. The degree theorem therefore gives exactly one preimage for every point between the boundary curves, and none outside. Local invertibility proves that $h_\varepsilon$ is a homeomorphism onto $\Omega_\varepsilon$. Each summand in \eqref{eq:689-map} is holomorphic or antiholomorphic, so the map is harmonic.

\paragraph{Second-order modulus formula.}
We record a short boundary-perturbation calculation used twice below. Work on the logarithmic cylinder with coordinate $s+i\phi$, of imaginary period $2\pi$. Suppose the boundaries of a perturbed ring are
\begin{equation}
 s=s_j+\delta p_j\cos(k\phi)
       +\delta^2(q_j+\text{zero-mean Fourier terms})+O(\delta^3),
 \quad s_0=0,\ s_1=m.
 \label{eq:689-generic-boundaries}
\end{equation}
The remainder is uniform with its first two derivatives. Let $A,B$ solve
\begin{equation}
 A+B=p_0,\qquad Ae^{km}+Be^{-km}=p_1,
 \qquad t_j=Ae^{ks_j}-Be^{-ks_j}.
 \label{eq:689-first-order-constants}
\end{equation}
Then its modulus is
\begin{equation}
 m+\delta^2\left[q_1-q_0-\frac{k}{2}(p_1t_1-p_0t_0)\right]+O(\delta^3).
 \label{eq:689-general-modulus}
\end{equation}
To verify this, use the periodic conformal approximation
\[
 W(\zeta)=\zeta+\delta(Ae^{k\zeta}+Be^{-k\zeta})+\delta^2K(\zeta).
\]
On the $j$th source boundary the angular variable becomes
$\phi=y+\delta t_j\sin(ky)+O(\delta^2)$. Substitution into
\eqref{eq:689-generic-boundaries} makes the mean of the real boundary coordinate
\[
 s_j+\delta^2\left(q_j-\frac{k}{2}p_jt_j\right)+O(\delta^3).
\]
The mean of $\Re K$ is the same on both straight boundaries. Hence their separation changes by the difference of these two second-order means, proving the coefficient in \eqref{eq:689-general-modulus}.

Here the calculation can be justified without assuming an unproved smooth dependence of uniformization. The second-order boundary discrepancies are finite Fourier sums. For each nonzero Fourier mode, two periodic exponential terms in $K$ match the two boundary discrepancies; the relevant two-by-two system is invertible since $m>0$. The two mean discrepancies determine the source-circle shifts. The resulting finite-Laurent map is injective for small $\delta$, since it is $C^1$-close to the identity on a slightly larger closed annulus. Its boundary mismatch is $O(\delta^3)$. Moving the source boundaries inward and outward by $C|\delta|^3$, for a sufficiently large fixed $C$, sandwiches the desired ring between two such conformal images. Modulus monotonicity bounds the error by $O(\delta^3)$. The same construction with finitely many real parameters gives a joint quadratic Taylor polynomial with a uniform cubic remainder: the boundary data and the Fourier correction coefficients depend polynomially on those parameters to this order.

\paragraph{Joint analytic dependence.}
There is also a fixed-domain justification of the uniform Taylor remainder used below. On $A_R=\{1<|w|<R\}$ the smooth map
\[
F_\varepsilon(w)=w+\varepsilon(2|w|^2-1)(w/|w|)^2
\]
maps both boundary circles exactly to \eqref{eq:689-boundaries} and is a diffeomorphism for small complex $\varepsilon$. For small complex $\beta$, the composition $L_\beta\circ F_\varepsilon$, where $L_\beta(z)=z+\beta\bar z$, is a real-analytic parameter family of diffeomorphisms. Pulling back the condenser energy gives a uniformly elliptic quadratic energy on the fixed annulus, with matrix coefficients real-analytic in the four real parameters. After fixing a smooth lift of the boundary values, Lax--Milgram and the Neumann series for the perturbed inverse give real-analytic dependence of the $H^1_0$ minimizer and its energy. The modulus is $2\pi$ divided by this positive capacity, and is therefore jointly real-analytic as well. This supplies the uniform cubic Taylor remainder independently of the finite-Laurent construction.

\paragraph{The harmonic perturbation decreases modulus.}
For a curve $z=r e^{i\theta}+\varepsilon c e^{2i\theta}$, put $d=c/r$. Expanding in the physical angular coordinate $\phi=\arg z$ yields
\begin{equation}
 \log|z|=\log r+\varepsilon d\cos\phi
    +\varepsilon^2\left(\frac{d^2}{2}-d^2\cos2\phi\right)
    +O(\varepsilon^3).
 \label{eq:689-harmonic-boundary-expansion}
\end{equation}
Indeed, $\phi=\theta+\varepsilon d\sin\theta+O(\varepsilon^2)$; substituting its inverse into $\Re\log(1+\varepsilon d e^{i\theta})$ gives the displayed formula. For the two curves in \eqref{eq:689-boundaries}, use $k=1$, $p_0=1$, $p_1=S/R$, and $q_j=p_j^2/2$ in \eqref{eq:689-general-modulus}. The constants in \eqref{eq:689-first-order-constants} are $A=2$ and $B=-1$. It follows that
\begin{equation}
 \operatorname{Mod}(\Omega_\varepsilon)
 =m-\left(1-R^{-2}\right)\varepsilon^2+O(\varepsilon^3).
 \label{eq:689-ordinary-modulus}
\end{equation}

\paragraph{Local affine optimization.}
After composition with a similarity and, if needed, conjugation, any invertible real affine map relevant to modulus is represented by
\[
 L_\beta(z)=z+\beta\bar z,\qquad |\beta|<1.
\]
Translations and similarities do not affect modulus. Set
\[
 M(\varepsilon,\beta)=\operatorname{Mod}(L_\beta\Omega_\varepsilon).
\]
First take $\varepsilon=0$ and $\beta=b$ real. On the image of a circle,
\begin{equation}
 \log|r(e^{i\theta}+b e^{-i\theta})|
 =\log r+b\cos2\phi
        +b^2\left(-1+\frac12\cos4\phi\right)+O(b^3).
 \label{eq:689-affine-boundary-expansion}
\end{equation}
Now \eqref{eq:689-general-modulus} applies with $k=2$, $p_0=p_1=1$, $q_0=q_1=-1$, and
\[
 A=\frac1{R^2+1},\qquad B=\frac{R^2}{R^2+1}.
\]
Consequently,
\begin{equation}
 M(0,b)=m-2\frac{R^2-1}{R^2+1}b^2+O(b^3).
 \label{eq:689-affine-modulus}
\end{equation}
Rotational invariance makes the complex-$\beta$ quadratic term a multiple of $|\beta|^2$ with the same coefficient.

There are no mixed quadratic terms between the harmonic and affine perturbations. To see this, temporarily allow $\varepsilon$ to be complex in the boundary parametrizations. Rotation by an angle $\alpha$ and boundary reparametrization give
\begin{equation}
 M(\varepsilon,\beta)=M(e^{-i\alpha}\varepsilon,e^{2i\alpha}\beta).
 \label{eq:689-rotation-identity}
\end{equation}
A quadratic monomial with one factor from $\varepsilon,\bar\varepsilon$ and one from $\beta,\bar\beta$ cannot be invariant under this action. Equivalently, the first-order perturbations have angular frequencies one and two, whose pairwise products have no constant Fourier coefficient. The joint Taylor construction above, together with \eqref{eq:689-ordinary-modulus} and \eqref{eq:689-affine-modulus}, therefore yields
\begin{equation}
 M(\varepsilon,\beta)
 =m-q|\varepsilon|^2-a|\beta|^2
       +O\bigl((|\varepsilon|+|\beta|)^3\bigr),
 \qquad
 q=1-R^{-2}>0,\quad a=2\frac{R^2-1}{R^2+1}>0.
 \label{eq:689-joint-expansion}
\end{equation}
In a fixed neighborhood of $(0,0)$ this gives
\begin{equation}
 M(\varepsilon,\beta)
 \le m-\frac q2|\varepsilon|^2-\frac a2|\beta|^2.
 \label{eq:689-local-upper}
\end{equation}

\paragraph{Excluding affine maps far from the identity.}
For a smooth ring whose two boundary curves enclose areas $A_0<A_1$, Carleman's area inequality and its equality case are
\begin{equation}
 \operatorname{Mod}(\Omega)\le\frac12\log(A_1/A_0),
 \quad\text{with equality only for a circular annulus up to similarity.}
 \label{eq:689-carleman}
\end{equation}
The needed version has a short proof. A conformal parametrization
$f(w)=\sum_{k\in\mathbb Z}a_kw^k$ from $1<|w|<e^M$, taking inner boundary to inner boundary, satisfies
\[
 A(r)=\pi\sum_{k\in\mathbb Z}k|a_k|^2r^{2k}.
\]
For analytic Jordan boundaries the map extends across the boundaries, so this area identity and its boundary limits are valid. Then
\[
 A(e^M)-e^{2M}A(1)
 =\pi\sum_k k|a_k|^2(e^{2kM}-e^{2M})\ge0.
\]
All nonzero summands are positive except those for $k=0,1$. Equality forces $f(w)=a_0+a_1w$, proving \eqref{eq:689-carleman} with its equality case.
For the image of the round annulus under $L_\beta$, the enclosed-area ratio remains $R^2$, since the affine determinant cancels. If $\beta\ne0$ the boundary curves are noncircular ellipses. Thus
\begin{equation}
 M(0,\beta)<m\qquad(\beta\ne0).
 \label{eq:689-strict-affine-max}
\end{equation}
On a compact set of parameters bounded away from both $0$ and the unit circle, this strict deficit is uniform and persists for sufficiently small $\varepsilon$ by continuity under smooth boundary perturbations. That continuity also follows from the fixed-domain Dirichlet-energy formulation.

It remains to control $|\beta|\uparrow1$ uniformly. The bounded inner complementary component of $\Omega_\varepsilon$ contains the disk of radius $1-|\varepsilon|$, while its outer boundary lies inside the disk of radius $R+|\varepsilon|S$. Hence, by ring-modulus monotonicity,
\begin{equation}
 M(\varepsilon,\beta)
 \le\operatorname{Mod}\left(L_\beta
        \{1-|\varepsilon|<|z|<R+|\varepsilon|S\}\right).
 \label{eq:689-ellipse-comparison}
\end{equation}
Up to rotation, $L_\beta$ has axes $1+|\beta|$ and $1-|\beta|$. For the elliptical ring obtained from $r<|z|<s$, every horizontal coordinate
$|x|\le(1+|\beta|)r/2$ determines an upper vertical segment joining the two boundary ellipses. Its length is at most $(1-|\beta|)s$. Every admissible condenser potential changes by one on this segment. Cauchy--Schwarz and integration in $x$ bound its Dirichlet energy below by
\[
 \frac{(1+|\beta|)r}{(1-|\beta|)s}.
\]
Since ring modulus is $2\pi$ divided by condenser capacity, the modulus is at most
\[
 2\pi\frac{(1-|\beta|)s}{(1+|\beta|)r}.
\]
For the radii in \eqref{eq:689-ellipse-comparison}, $r$ stays bounded below and $s$ bounded above as $\varepsilon\to0$. Thus $M(\varepsilon,\beta)\to0$ uniformly as $|\beta|\uparrow1$.

Choose a small parameter neighborhood in which \eqref{eq:689-local-upper} holds. The compact strict bound from \eqref{eq:689-strict-affine-max} and the degeneration estimate show that all remaining affine parameters have modulus at most $m-\gamma$ for some fixed $\gamma>0$, for sufficiently small $|\varepsilon|$. Shrinking $|\varepsilon|$ further if necessary, the full affine supremum is bounded by
\[
 \operatorname{Mod}_{\mathrm{aff}}(\Omega_\varepsilon)
 \le m-\frac q2\varepsilon^2<m
 \qquad(\varepsilon\ne0).
\]
Together with the harmonic homeomorphism \eqref{eq:689-map}, this proves the theorem.
\end{proof}

\paragraph{Independent numerical check.}
The proof is asymptotic and does not depend on numerical optimization. As a cross-check, at $\tau=1$ the baseline is $m=0.4337808304830271$. Harmonic Laurent collocation gives
$(\operatorname{Mod}(\Omega_\varepsilon)-m)/\varepsilon^2$ equal to
$-0.5841929$, $-0.5810615$, $-0.5802843$, and $-0.5800670$ at
$\varepsilon=0.02,0.01,0.005,0.002$, respectively, approaching the proved coefficient $-(1-R^{-2})=-0.5800256584$. At $\varepsilon=0.02$, numerical affine optimization gives $\beta\simeq0.003933496$ and modulus $0.433559780729$, still below the proposed bound. These values provide a consistency check only.

\begingroup\small
\begin{thebibliography}{99}
\bibitem{IwaniecKovalevOnninen} T. Iwaniec, L. V. Kovalev and J. Onninen, \emph{Harmonic mapping problem and affine capacity}, Proc. Roy. Soc. Edinburgh Sect. A \textbf{141} (2011), 1017--1030. \url{https://arxiv.org/abs/1001.2124}.
\end{thebibliography}
\endgroup
\end{document}
